\section{방법}
\label{sec:method}

\begin{figure*}[t]
  \centering
  \includegraphics[width=\linewidth]{figures/fig_overview.pdf}
  \caption{학습에 쓰지 않은(held-out) Grand Central 상황에서의 개요(로봇: 위쪽을 향하는 검은 원; 사람: 속도 표시가 있는 회색 원;
  녹색 점선 화살표: 목적지 방향이며, 목적지는 화면 밖에 있다).
  (a) VQ-VAE 코드북은 로봇의 상태로부터 $K=16$개의 후보 경로를 디코딩한다.
  (b) 안전 필터는 등속 예측 하에서 모든 사람과 최소 0.5\,m 거리를 유지하고 진행량이 있는 경로만 남긴다.
  (c) 이 상황에서 EUBO가 던질 질문: 더 나은 선택지의 기대 효용이 가장 높은 쌍.
  (d) 다섯 개의 스타일 특징에 대한 베이지안 Bradley--Terry 모델이 선호 $\hat\theta$를 추정한다. 여기서는 사람 뒤로 지나가는 것과 작은 우회를 선호하는
  예시 사용자 $\theta^\ast$가 다른 다섯 상황에서 다섯 번 답한 후의 추정이다.
  (e) 그 다섯 상황에 포함되지 않은 이 상황에서, 로봇은 $\hat\theta$ 하에서 가장 좋은 안전한 경로를 실행한다. 이 경로는 가장 사람다운 경로
  (인간 사전분포(human prior))와 다르다. \texttt{utils/figure\_overview.py}로 생성하였다.}
  \label{fig:overview}
\end{figure*}

\subsection{문제 설정}
상황 $s$에서 로봇은 자신의 상태(속력, 과거 두 위치), 목적지, 그리고 주변 사람들의 위치와 속도를 알고 있다.
사용자는 로봇이 걷는 방식에 대한 선호를 가지며, 우리는 이를 스타일 특징에 대한 선형 효용으로 모델링한다.
사용자는 몇 개의 질문 상황에서 소수의 A/B 질문(query)에 답하고, 이후 로봇은 사용자가 한 번도 질문받지 않은 새로운 상황에서
자신의 경로를 선택해야 한다.

\subsection{보행 행동 코드북}
\label{sec:candidates}
$K=16$개의 코드를 갖는 VQ-VAE~\citep{vandenoord2017neural}를, 진행 방향 좌표계에서 등속으로부터의 편차로 표현한
보행자 미래 궤적(0.8\,s 간격 6개 점, 4.8\,s)으로 학습한다. 로봇의 현재 속력에서 모든 코드를 디코딩하면 $K$개의 후보 경로를 얻는다.
같은 데이터로 목적지와 주변 사람 요약을 입력으로 하여 학습한 사전 네트워크 $p(\text{code} \mid s)$는
각 후보가 그 상황에 있는 사람에게 얼마나 전형적인지를 점수화한다.
코드북과 사전분포는 데이터셋마다 하나씩, 사람 단위 분할(약 80/10/10\%, 어떤 사람도 두 분할에 동시에 속하지 않음)의 학습 분할로 학습한다.
모든 평가 상황은 테스트 분할에서 가져온다. 학습 샘플은 $\ge$0.5\,m/s로 걷고 있으며 미래 주석 6개를 사용할 수 있는 임의의 주석된 보행자와,
6\,m 이내의 최대 12명의 사람으로 구성된다. VQ-VAE($K=16$개 코드, 8차원 잠재 변수)는 Adam(학습률 $10^{-3}$)으로
Grand Central에서는 30 에폭(배치 1024), 더 작은 ETH/UCY에서는 150 에폭(배치 256) 동안 학습한다.
사전분포는 256개 유닛의 은닉층 두 개를 갖는 MLP이며, Adam($10^{-3}$, weight decay $10^{-5}$)으로 30 에폭 동안,
목적지가 의미를 가질 만큼 충분히 앞에 있는 샘플(Grand Central에서는 12\,s, ETH/UCY에서는 8\,s)로 학습한다.
입력은 속력, 과거 두 위치, 목적지(트랙의 마지막 위치), 그리고 주변 사람에 대한 8개 수치 요약
(앞쪽의 가장 가까운 사람, $<$1\,m 접근까지의 최소 시간, 앞쪽 좌측과 우측의 사람)이다. 데이터셋 크기는 Table~\ref{tab:data}에 있다.

\paragraph{안전 필터.}
사람들은 등속으로 외삽한다. 후보는 모든 사람과의 최소 거리가 $c=0.5$\,m 이상이고, 목적지를 향한
진행량이 가장 좋은 후보의 $\rho=0.5$배 이상이면 안전하다. 선호는 안전한 후보들 중에서만 선택한다:
\begin{equation}
  \mathcal{C}(s) = \bigl\{ k \;:\; d_{\min}(s,k) \ge c,\; \mathrm{prog}(s,k) \ge \rho \max_{k'} \mathrm{prog}(s,k') \bigr\}.
  \label{eq:safe}
\end{equation}
안전한 후보가 네 개 미만인 상황은 사용하지 않는다.

\subsection{스타일 특징}
각 안전한 후보는 $\phi(s,k) \in \mathbb{R}^5$로 기술되며, 어느 특징도 사전에 정해진 올바른 부호를 갖지 않는다:
\emph{side}(최근접 시점에서 충돌 대상 사람으로부터의 횡방향 오프셋의 $\tanh$),
\emph{front}(그 사람의 앞으로 가로지르는지 뒤로 가로지르는지), \emph{yield}(현재 속력 대비 평균 속력),
\emph{detour}(목적지까지의 직선으로부터의 최대 편차), \emph{human}($\log p(\text{code}\mid s)$).
특징은 질문 상황들의 안전한 후보들에 대해 표준화한다.

\subsection{선호 학습과 질문 선택}
\label{sec:preference}
사용자의 효용은 $u(s,k) = \theta^\top \phi(s,k)$이다. 응답은 Bradley--Terry 모델~\citep{bradley1952rank}
$P(i \succ j \mid \theta) = \sigma(\theta^\top(\phi_i - \phi_j))$을 따르며 사전분포는 $\theta \sim \mathcal{N}(0, I)$이다.
사후분포는 라플라스 근사(MAP에는 뉴턴법, 공분산에는 역 헤시안)로 근사한다.
각 질문은 최선 선택지의 기대 효용을 최대화하는 쌍을 묻는다~\citep{astudillo2023qeubo}:
\begin{equation}
  (i,j) = \operatorname*{arg\,max}_{i \ne j} \; \mathbb{E}_{\theta \sim p(\theta \mid \mathcal{D})}\bigl[\max\{u_i, u_j\}\bigr],
\end{equation}
이는 512개의 사후분포 샘플로 추정한다. 테스트 시에는 어떤 상황에서든 로봇은 다음을 실행한다.
\begin{equation}
  \hat k_n(s) = \operatorname*{arg\,max}_{k \in \mathcal{C}(s)} \hat\theta_n^{\top}\phi(s,k),
  \label{eq:exec}
\end{equation}
여기서 $\hat\theta_n$은 $n$개의 응답 후의 사후분포 평균이다.

% Formal regret definition (11/18 feedback item 4). Setting, safe set and executed path are defined in Sec. 3.
\subsection{평가: 단순 regret}
\label{sec:regret}

\paragraph{사용자.}
평가에서 사용자는 효용 $u^\ast(s,k) = \theta^{\ast\top}\phi(s,k)$를 갖는 단위 벡터 $\theta^\ast$($\lVert\theta^\ast\rVert = 1$)이며,
질문(query) $(i,j)$에 $P(i \succ j) = \sigma\bigl(\beta\,(u^\ast(s,i) - u^\ast(s,j))\bigr)$, $\beta = 5$의 확률로 답한다.
$n$개의 응답 후 방법은 $\hat\theta_n$(Bradley--Terry 방법의 경우 라플라스 사후분포 평균, DPO의 경우 DPO 해,
Sec.~\ref{sec:dpo})을 반환하고, Eq.~\eqref{eq:safe}의 안전한 후보 $\mathcal{C}(s)$ 중에서 Eq.~\eqref{eq:exec}의 $\hat k_n(s)$를 실행한다.

\paragraph{단순 regret.}
그 선택의 regret은 사용자에게 가장 좋은 안전한 후보 대비 잃은 효용이며,
상황 간 비교가 가능하도록 해당 상황의 효용 범위로 정규화한다:
\begin{equation}
  r_n(s) = \frac{\max_{k \in \mathcal{C}(s)} u^\ast(s,k) \;-\; u^\ast\bigl(s,\hat k_n(s)\bigr)}
                {\max_{k \in \mathcal{C}(s)} u^\ast(s,k) \;-\; \min_{k \in \mathcal{C}(s)} u^\ast(s,k)}
  \;\in [0,1].
  \label{eq:regret}
\end{equation}
보고하는 수치는 사용자와 학습에 쓰지 않은(held-out) 테스트 상황에 대한 이 값의 기댓값이다:
\begin{equation}
  \mathrm{SR}_n = \mathbb{E}_{\theta^\ast}\Bigl[\frac{1}{|\mathcal{S}_{\text{test}}|}
  \sum_{s \in \mathcal{S}_{\text{test}}} r_n(s)\Bigr],
  \label{eq:sr}
\end{equation}
여기서 응답은 테스트 상황과 겹치지 않는 질문 상황에서만 수집한다
($\mathcal{S}_{\text{query}} \cap \mathcal{S}_{\text{test}} = \emptyset$; 크기는 Sec.~\ref{sec:experiments} 참조).
오라클($\hat\theta = \theta^\ast$)의 경우 $\mathrm{SR}_n = 0$이며,
안전한 후보를 균등하게 무작위로 고르면 $\mathrm{SR} \approx 0.5$가 된다.

\paragraph{왜 단순 regret인가.}
질문은 사용자의 주의를 소모하지만 주행 효용은 소모하지 않으며, 중요한 것은 로봇이 그 이후에
실행하는 경로이다. 따라서 목표는 질문에 걸친 누적 regret이 아니라 최선 팔 식별(best-arm identification; 단순 regret(simple regret))이다.
EUBO~\citep{astudillo2023qeubo}는 잡음 없는 응답 하에서 바로 이 목표에 대해 한 스텝 베이즈 최적이며,
이것이 EUBO를 사용하는 이유이다.
Fig.~\ref{fig:regret}은 질문 수 $n \in \{0,2,5,10,20\}$에 따른 $\mathrm{SR}_n$을 나타낸다.

\paragraph{다른 곳에서 사용하는 변형.}
\emph{폐루프:} 주행 중 0.8\,s마다 이루어지는 각 재계획 스텝 $t$에서, Eq.~\eqref{eq:regret}을
그 스텝의 안전 집합 $\mathcal{C}(s_t)$에 대해 평가하고 에피소드의 스텝들에 걸쳐 평균한다(스텝 regret).
\emph{후보 생성기:} 같은 상황에서 두 생성기 $A$와 $B$를 비교하기 위해, Eq.~\eqref{eq:regret}의 max와 min을
$\mathcal{C}_A(s) \cup \mathcal{C}_B(s)$에 대해 취한다(합집합 regret). 따라서 후보들이 사용자가 원하는 것을 놓치는 생성기는
자신의 후보 중 최선을 고르더라도 불이익을 받는다.
생성기별로 계산한 Eq.~\eqref{eq:regret}의 regret은 생성기 간에 비교할 수 없다.
\emph{파일럿:} 실제 참가자의 경우 $\theta^\ast$를 알 수 없으므로, 측정 가능한 대리 지표는
학습에 쓰지 않은(테스트 블록) 응답 중 $\hat\theta$가 예측하는 비율이다.

% Figure (utils/plot_regret_curves.py -> reports/fig_regret.pdf)
\begin{figure}[t]
  \centering
  \includegraphics[width=\linewidth]{figures/fig_regret.pdf}
  \caption{학습에 쓰지 않은 상황에서 A/B 질문 수에 따른 단순 regret $\mathrm{SR}_n$(Eq.~\ref{eq:sr}).
  음영: 200명의 합성 사용자에 대한 95\% 부트스트랩 CI. 점선: 인간 사전분포(human prior)만 사용; 오라클은 0이다. BT: Bradley--Terry(라플라스); DPO: Sec.~\ref{sec:dpo}.}
  \label{fig:regret}
\end{figure}

